\chapter{Kidney Stone Analysis}
\section*{What Is This Test?} Kidney stone analysis identifies the chemical composition of a passed or removed urinary stone.\section*{Alternative Names} Urinary calculus analysis; renal stone analysis; nephrolithiasis composition.\section*{Principle} Infrared spectroscopy or X-ray diffraction identifies mineral components such as calcium oxalate, uric acid, struvite, or cystine.\section*{Why This Test Is Ordered} It helps explain stone formation and select prevention strategies after a stone is passed or removed.\section*{Primary vs Secondary Test} Stone composition is secondary; imaging, urinalysis, and metabolic urine studies assess the broader risk.\section*{How This Test Is Performed} The stone is collected in a clean container, dried as instructed, and sent to the laboratory.\section*{What Preparation Is Needed} Strain urine if instructed and avoid contaminating the specimen with tissue or toilet water.\section*{Understanding Results} Mixed or single mineral composition guides dietary and medical prevention but does not alone identify the cause.\section*{Follow-up Tests} Serum chemistries, urinalysis, and 24-hour urine calcium, oxalate, citrate, uric acid, sodium, and volume may follow.\section*{References} \begin{enumerate}\item MedlinePlus. Kidney Stone Analysis.\item American Urological Association. Medical management of kidney stones.\end{enumerate}\clearpage\section*{Understanding Results and Follow-up}
The laboratory report should identify the specimen, method, units, reference interval or decision limit, and whether the result is positive, negative, abnormal, or inconclusive. Reference intervals vary with age, sex, pregnancy, specimen type, assay, and laboratory; a result outside a range is not automatically a diagnosis, and a result within a range does not always exclude disease. Discuss unexpected results with the ordering clinician, who can decide whether repeat or confirmatory testing, treatment, or referral is appropriate. Do not change medicines or delay urgent care based on an isolated result.
\section*{Regulatory Status and Authorization Date}This chapter describes a test category, not a specific commercial product. FDA, CDSCO, EU IVDR, and MHRA approval or registration dates therefore cannot be assigned without the assay manufacturer, product name, instrument, and intended use. For laboratory-developed tests, an individual agency approval date may not exist. See the Preface for the interpretation of regulatory status.
\clearpage
